\chapter{ارزیابی، نتیجه‌گیری و کارهای آینده}

\section{روش ارزیابی}
برای ارزیابی سامانه، ماژول \lr{\texttt{profiler-evaluation}} سناریوهای کنترل‌شده‌ای را اجرا می‌کند. این ماژول یک فرایند جاوا راه‌اندازی می‌کند، شناسه فرایند را چاپ می‌کند و پس از یک تأخیر اولیه، سناریوها را اجرا می‌کند تا دریافت‌کننده‌های \lr{JFR} و هسته فرصت اتصال داشته باشند. خروجی ارزیابی شامل فایل‌های \lr{\texttt{experiment\_results.csv}}، \lr{\texttt{experiment\_metrics.json}} و \lr{\texttt{evaluation\_report.md}} است.

\section{سناریوهای آزمایش}
سناریوهای پیاده‌سازی‌شده در پروژه عبارت‌اند از:
\begin{itemize}
\item \lr{\texttt{CpuSpikeScenario}} برای ایجاد افزایش ناگهانی مصرف پردازنده؛
\item \lr{\texttt{MemoryGcPressureScenario}} برای ایجاد فشار حافظه و تحریک جمع‌آوری زباله؛
\item \lr{\texttt{FileIoBurstScenario}} برای تولید الگوی شدید ورودی/خروجی فایل؛
\item \lr{\texttt{LockContentionScenario}} برای ایجاد رقابت روی قفل‌ها؛
\item \lr{\texttt{NetworkBurstScenario}} برای ایجاد ترافیک شبکه؛
\item \lr{\texttt{MmapFileScanScenario}} برای بررسی رفتار دسترسی حافظه‌نگاشته به فایل.
\end{itemize}

هر سناریو یک رفتار هدفمند ایجاد می‌کند و سپس گزارش تحلیلی بررسی می‌شود تا مشخص شود آیا ناهنجاری مورد انتظار تشخیص داده شده است یا خیر. این روش برای پروژه حاضر مناسب است، زیرا امکان کنترل زمان شروع، مدت اجرا و نوع فشار ایجادشده را فراهم می‌کند.

\section{معیارهای ارزیابی}
دو معیار اصلی برای ارزیابی در نظر گرفته می‌شود. معیار نخست موفقیت تشخیص است؛ یعنی آیا سامانه برای سناریوی اجراشده، ناهنجاری مناسب را گزارش کرده است یا نه. معیار دوم تأخیر تشخیص است؛ یعنی فاصله زمانی بین شروع سناریو و ثبت ناهنجاری در گزارش. این دو معیار در کنار هم نشان می‌دهند سامانه هم از نظر دقت عملی و هم از نظر سرعت واکنش چه وضعیتی دارد.

\section{بحث}
طراحی خط لوله‌ای پروژه چند مزیت مهم دارد. نخست، هر جزء به صورت مستقل اجرا و عیب‌یابی می‌شود. دوم، استفاده از \lr{Kafka} امکان جداسازی نرخ تولید و مصرف پیام را فراهم می‌کند. سوم، تجمیع پنجره‌ای حجم داده خام را کاهش داده و تحلیل را ساده‌تر می‌کند. چهارم، ترکیب رخدادهای \lr{JFR} و \lr{eBPF} دید گسترده‌تری نسبت به رفتار برنامه ایجاد می‌کند.

در مقابل، سامانه محدودیت‌هایی نیز دارد. کیفیت تشخیص به انتخاب آستانه‌ها و طول پنجره وابسته است. همچنین اجرای بخش \lr{eBPF} به محیط لینوکس و دسترسی‌های لازم نیاز دارد. از سوی دیگر، بسیاری از آزمون‌های موجود در پروژه حالت اسکلت اولیه دارند و برای استفاده پژوهشی کامل، لازم است پوشش آزمون واحد و آزمون یکپارچه افزایش یابد.

\section{نتیجه‌گیری}
در این پایان‌نامه، سامانه‌ای برای پایش عملکرد و تشخیص ناهنجاری در برنامه‌های جاوا بررسی، طراحی و مستندسازی شد. سامانه با استفاده از \lr{JFR} رخدادهای سطح ماشین مجازی را دریافت می‌کند، با \lr{eBPF} امکان مشاهده رخدادهای سطح هسته را فراهم می‌سازد، با \lr{Kafka Streams} داده‌ها را در پنجره‌های زمانی تجمیع می‌کند و با سرویس تحلیلی پایتونی ناهنجاری‌ها را گزارش می‌دهد. نتیجه اصلی این پژوهش آن است که ترکیب مشاهده‌پذیری سطح برنامه و سطح سیستم‌عامل در یک معماری جریانی، امکان تحلیل دقیق‌تر و توسعه‌پذیرتر رفتار عملکردی را فراهم می‌کند.

\section{کارهای آینده}
برای توسعه آینده پروژه، موارد زیر پیشنهاد می‌شود:
\begin{itemize}
\item افزودن داشبورد تصویری برای نمایش زنده معیارها و ناهنجاری‌ها؛
\item استفاده از روش‌های یادگیری ماشین برای تنظیم پویا و داده‌محور آستانه‌ها؛
\item افزایش پوشش آزمون‌های واحد و یکپارچه برای ماژول‌های جاوا و پایتون؛
\item افزودن همبستگی بین رخدادهای \lr{JFR} و رخدادهای هسته بر اساس زمان و شناسه فرایند؛
\item پشتیبانی از استقرار کانتینری کامل برای اجرای ساده‌تر کل خط لوله؛
\item افزودن معیارهای دقت، فراخوانی و نرخ هشدار کاذب به گزارش ارزیابی.
\end{itemize}
